\documentclass[oneside, 11pt]{article}

\usepackage[T1]{fontenc}
\usepackage[utf8]{inputenc}
\usepackage[english]{babel}
\usepackage{enumerate}
\usepackage{isotope}


\usepackage{fouriernc}
\usepackage[detect-all, binary-units, separate-uncertainty=true,
            per-mode=symbol, retain-explicit-plus, retain-unity-mantissa=false]{siunitx}

\usepackage{setspace}
\setstretch{1.2}

\setlength{\parskip}{\smallskipamount}
\setlength{\parindent}{0pt}

\usepackage[headheight=14pt]{geometry}
\geometry{marginparwidth=0.5cm, verbose, a4paper, tmargin=3cm, bmargin=3cm,
          lmargin=2cm, rmargin=2cm}

\usepackage{float}

\usepackage[fleqn]{amsmath}
\numberwithin{equation}{section}
\numberwithin{figure}{section}

\usepackage{graphicx}
\graphicspath{{images/}{../../../images/}}

\usepackage{tikz}
\usetikzlibrary{shapes}
\usetikzlibrary{plotmarks}

\newcounter{Exercise}
\setcounter{Exercise}{1}
\usepackage{xcolor}
\definecolor{shadecolor}{gray}{0.9}
\usepackage{framed}
\usepackage{caption}

\usepackage{url}


\usepackage{fancyhdr}
\pagestyle{fancy}
\fancyhf{}
\rhead{\thepage}
\renewcommand{\footrulewidth}{0pt}
\renewcommand{\headrulewidth}{0pt}

\fancypagestyle{firststyle}
{
    \fancyhf{}
    \rhead{\thepage}
    \cfoot{\includegraphics[height=30pt]{HiSPARClogo}}
    \rfoot{\includegraphics[height=25pt]{CCbysa}}
    \lfoot{\includegraphics[height=30pt]{NIKHEFlogo}}
    \renewcommand{\footskip}{50pt}
    \renewcommand{\footrulewidth}{0.1pt}
    \renewcommand{\headrulewidth}{0pt}
}

\newcommand{\figref}[1]{Figuur~\ref{#1}}

\newcommand{\hisparc}{\textsmaller{HiSPARC}\xspace}
\newcommand{\kascade}{\textsmaller{KASCADE}\xspace}
\newcommand{\sapphire}{\textsmaller{SAPPHiRE}\xspace}
\newcommand{\jsparc}{\textsmaller{jSparc}\xspace}
\newcommand{\hdf}{\textsmaller{HDF5}\xspace}
\newcommand{\aires}{\textsmaller{AIRES}\xspace}
\newcommand{\csv}{\textsmaller{CSV}\xspace}
\newcommand{\python}{\textsmaller{PYTHON}\xspace}
\newcommand{\corsika}{\textsmaller{CORSIKA}\xspace}
\newcommand{\labview}{\textsmaller{LabVIEW}\xspace}
\newcommand{\daq}{\textsmaller{DAQ}\xspace}
\newcommand{\adc}{\textsmaller{ADC}\xspace}
\newcommand{\hi}{\textsc{h i}\xspace}
\newcommand{\hii}{\textsc{h ii}\xspace}
\newcommand{\mip}{\textsmaller{MIP}\xspace}
\newcommand{\hisparcii}{\textsmaller{HiSPARC II}\xspace}
\newcommand{\hisparciii}{\textsmaller{HiSPARC III}\xspace}

\DeclareSIUnit{\electronvolt}{\ensuremath{\mathrm{e\!\!\:V}}}

\DeclareSIUnit{\unitsigma}{\ensuremath{\sigma}}
\DeclareSIUnit{\mip}{\textsmaller{MIP}}
\DeclareSIUnit{\adc}{\textsmaller{ADC}}

\DeclareSIUnit{\gauss}{G}
\DeclareSIUnit{\parsec}{pc}
\DeclareSIUnit{\year}{yr}



\begin{document}

\title{Tsjerenkov}
\author{N.G. Schultheiss}
\date{}

\maketitle
\thispagestyle{firststyle}

\section{Inleiding}

Deze module volgt op de module ``Compton'' en behandelt de Tsjerenkov-verstrooiing.
Er is sprake van Tsjerenkov-verstrooiing als een deeltje door een
stof beweegt met een snelheid die groter is dan de lichtsnelheid in
de stof. De snelheid van het licht in een stof is te berekenen met
het Debye-Einstein model. Omdat we kijken naar ionen met een tegengestelde
lading zullen ze tegen elkaar in bewegen door het elektromagnetisch
veld. Volgens de Debye-Einstein is er overigens ook een acoustische
golf mogelijk. De ionen bewegen nu min of meer synchroon. Met deze
acoustische golven is de soortelijke warmte van stoffen voor een deel
te verklaren. Bij geleidende stoffen zit een deel van de soortelijke
warmte in het elektronen gas. De soortelijke warmte van stoffen valt
echter buiten dit betoog. 


\section{Tsjerenkov}

Pavel Aleksejevitsj Tsjerenkov onderzocht de straling die bijvoorbeeld
in het reactorvat van een kerncentrale plaatsvindt. Dit vat is meestal
gevuld met water. Water werkt in dit geval zowel als koelmiddel als als
moderator\footnote{Een moderator remt de neutronen bij een
kernreactie af, zodat ze weer kunnen reageren met een volgende te
splitsen kern. Zonder moderator zullen de te snelle neutronen uit de
reactor ontsnappen zonder dat ze een nieuwe reactie starten. De reactor
stopt dan. Bij teveel trage neutronen loopt de reactie ook uit de hand.
Teveel kernen reageren en we spreken van een melt-down (Harrisburg /
Tsjernobyl).}.

Zoals we hebben gezien bij de Geigerteller kunnen deeltjes electronen
in een medium wegbotsen. De electronen springen dan als het ware naar
een een hogere aangeslagen energietoestand. Bij Tsjerenkov-straling
vallen deze electronen praktisch direct terug en zenden de energie
als foton uit. Als dit in water gebeurt, zal de snelheid van het foton
niet de lichtsnelheid zijn maar $c_{medium}=\frac{\mathrm{c}}{N_{medium}}$.
Het deeltje kan echter wel een hogere snelheid dan $c_{medium}$ hebben.
We nemen nu een deeltje waar dat sneller gaat dan de lichtsnelheid
in het medium. Net als bij een supersoon vliegtuig zendt het deeltje
een kegel licht (in plaats van geluid) uit. Omdat het deeltje door
het botsen met de elektronen ook een verstrooiing krijgt, is deze
kegel overigens niet scherp begrensd. Dit verschijnsel wordt Tsjerenkov-straling
genoemd. Lichte deeltjes zoals electronen hebben een grotere verstrooiing
dan zwaardere deeltjes. Als het licht van de kegel wordt gedetecteerd
zal eer voor zware deeltjes een korte puls worden wargenomen. Omdat
de kegel bij lichtere deeltjes heen en weer beweegt, wordt de gemeten
puls iets breder.

\begin{figure}[H]
\noindent \begin{centering}
\includegraphics[scale=0.75]{Tsjerenkov-cone}
\par\end{centering}

\caption{Een schematische voorstelling van Tsjerenkov-straling}
\end{figure}


\section{Licht in stoffen}

Als een elektromagnetische golf op een stof valt, zullen de ionen
in de vaste stof door het wisselend elektrisch veld heen en weer bewegen.
Deze beweging gaat als een golf door de stof. Op een gegeven moment
heeft de golf het einde van de stof bereikt. Er schudden nu ionen
heen en weer, deze veroorzaken weer een elektromagnetische golf.

\begin{figure}[H]
\noindent \begin{centering}
\includegraphics[scale=0.75]{rooster1}
\par\end{centering}

\caption{Licht in stoffen}
\end{figure}


Als er een geladen deeltje door de stof vliegt, zullen de deeltjes
ook heen en weer bewegen. Deze beweging verplaatst zich als een golf
door de stof. Aan de overgang van stof naar vacuum wordt de beweging
van de ionen omgezet in een elektromagnetische golf.

\begin{figure}[H]
\noindent \begin{centering}
\includegraphics[scale=0.75]{rooster2}
\par\end{centering}

\caption{Een geladen deeltje in een rooster}
\end{figure}

\end{document}
